\chapter{Prova}

Qui si inserisce il contenuto; è comodo mettere
\begin{verbatim}
    \chapter{titolo del capitolo}
    \input{file_contenente_il_testo_del_capitolo.tex}
\end{verbatim}

così da gestire gestire separatamente i vari capitoli e non avere un file unico troppo lungo da gestire.

Tutte le immagini devono essere messe nella cartella \texttt{images} (o in quella indicata come \texttt{graphicspath} nel file \texttt{main.tex}) e poi possono essere inserite comodamente nel testo con il comando: 
\begin{verbatim}
    \includegraphics[width=.5\textwidth]{nome_file.png}
\end{verbatim}

Per le citazioni è sufficiente fare \cite{felix20}, inserendo come argomento del comando \verb|\cite| lo stesso nome della citazione che è stato usato nel file \texttt{bib}. Se non si inserisce alcuna citaziona nel testo col comando \verb|\cite| non verrà generata alcuna bibliografia.